\documentclass[10pt,a4paper]{amsart}
\usepackage[margin=19mm]{geometry}
\usepackage{amsmath,amssymb,mathtools}
\usepackage[scr=rsfs,cal=euler]{mathalfa}
\usepackage{xfrac}
\usepackage[unicode,hidelinks,bookmarks=false]{hyperref}
\newcommand{\ie}{i.e.\ }
\newcommand{\cf}{cf.\ }
\newcommand{\resp}{respectively\ }
\newcommand{\Subsection}{\subsection}
\providecommand{\nobreakdash}{\nobreak\hbox{-}}
\setlength{\emergencystretch}{2em}
\multlinegap0pt
\usepackage{bm}
\usepackage{bbm}
\usepackage{dsfont}
\usepackage{xspace}
\usepackage{cancel}
\usepackage{mathtools}
\usepackage{amsthm}
\makeatletter
\@ifundefined{thm}{\newtheorem{thm}{Thm}}{}
\@ifundefined{lem}{\newtheorem{lem}[thm]{Lemma}}{}
\@ifundefined{prop}{\newtheorem{prop}[thm]{Proposition}}{}
\makeatother
\DeclareMathOperator{\Par}{\mathbb{Y}}
\title[Higher Order Bell Symmetric Functions]{Higher Order Bell Symmetric Functions}
\author{Milo Bechtloff Weising}
\date{Source: arXiv:2505.18504v2 (CC BY 4.0)}
\begin{document}
\maketitle

\noindent\emph{Source and licence.} Higher Order Bell Symmetric Functions. Version arXiv:2505.18504v2 (CC BY 4.0), distributed under the Creative Commons Attribution 4.0 International licence, as recorded in the arXiv licence metadata for this exact version and shown on the abstract page https://arxiv.org/abs/2505.18504v2. Full rights notice: licenses/arxiv-2505.18504-CC-BY-4.0.txt.\par\smallskip

\noindent\textit{Editorial scope.} Counted unit: the Prop whose source label is \texttt{None} in the article (source0.tex lines 645--648, source label \texttt{None}), delivered together with the article's own complete original proof of it (source0.tex lines 649--675, one \texttt{proof} environment of 27 lines). No mathematics is written, repaired or re-derived: statement and proof are the article's own text line for line. Required context retained uncounted: monomial recursion prop (lines 367--369); restriction lemma (lines 622--624). Every definition, display and label of those lines is retained unchanged. Omitted from this package and not redistributed: 1--366, 370--621, 625--644, 676--745. Nothing inside a retained excerpt is omitted or altered: every retained body is delivered exactly as the article's lines are written, so no omission receipt of any kind accompanies this package. Citations: the retained excerpts carry no citation command at all, so no bibliography entry of the article is transcribed and no reference is redistributed. Reference adaptation: none was needed and none was made; every reference inside the delivered unit resolves inside it, so no unresolved reference or citation is produced. Printed slips kept literal: no printed word, symbol, hypothesis or number of the counted statement or of its proof was altered. Layout and class: the article's own class is \texttt{amsart} with packages including geometry, amssymb, amsmath, mathrsfs, amsfonts, hyperref, xcolor, bbm, tikz, array, multirow, ytableau, thmtools, thm-restate; the wrapper is the lane's pinned \texttt{amsart} editorial wrapper at 10pt on A4, which restates the theorem-like environments the retained text needs and adds the article's own macro definitions that the retained text needs (\texttt{\textbackslash Par}), with extra wrapper packages (bm, bbm, dsfont, xspace, cancel, mathtools). The delivered numbering is the unit's own. Assets: no retained span contains a figure environment, a graphics-inclusion command, a PDF-page inclusion, a table or a TikZ picture; no image file is copied into this package, the package declares no asset dependency and no image file is redistributed with it. Licence: version arXiv:2505.18504v2 of this article is distributed under the Creative Commons Attribution 4.0 International licence, as recorded in the arXiv licence metadata for this exact version and shown on the abstract page https://arxiv.org/abs/2505.18504v2. A full rights notice is delivered at licenses/arxiv-2505.18504-CC-BY-4.0.txt. Characters: every retained excerpt is pure ASCII, so the delivered engine, XeLaTeX with its default Latin Modern text font, reports no missing character.
\begin{prop}\label{monomial recursion prop} For all $m \geq 0$ and $\lambda \in \Par \setminus \{ \emptyset \},$
    $$\rho^{(m+1)}(\lambda) = \sum_{\substack{\phi:Q'\rightarrow \mathbb{Z}_{\geq 0} \\ \sum_{\alpha \in Q'} \phi(\alpha)\alpha = \lambda}} \prod_{\alpha \in Q'} \binom{\phi(\alpha) + \rho^{(m)}(\alpha)-1}{\rho^{(m)}(\alpha)-1}.$$
\end{prop}

\begin{lem}\label{restriction lemma}
    $$\Omega\left(z \Omega(X) \right) = \sum_{\lambda} R_{\lambda}(z) s_{\lambda}(X)$$
\end{lem}

\begin{prop}
    For all $\lambda \in \Par,$
     $(1-z)R_{\lambda}(z)$ is a polynomial with degree at most $|\lambda|$. 
\end{prop}

\begin{proof}
Using Lemma \ref{restriction lemma} we see that 
$$R_{\lambda}(z) = \langle s_{\lambda}(X), \Omega\left(z \Omega(X) \right) \rangle.$$ We will use the expansion $s_{\lambda} = \sum_{\mu \vdash n} K^{(-1)}_{\lambda,\mu} h_{\mu}$ where $K^{(-1)}_{\lambda,\mu}$ are the \textbf{inverse Kostka numbers} and instead compute $\langle h_{\mu}(X), \Omega\left(z \Omega(X) \right) \rangle$. Following similar to the proof of Proposition \ref{monomial recursion prop}, 

\begin{align*}
    &\Omega\left(z\Omega(X) \right) \\
    &= \Omega\left( \sum_{\alpha \in Q} zx^{\alpha} \right)\\
    &= \prod_{\alpha \in Q} \frac{1}{1-zx^{\alpha}}\\
    &= \frac{1}{1-z}\prod_{\alpha \in Q\setminus \{0\}} \frac{1}{1-zx^{\alpha}}\\
    &= \frac{1}{1-z} \sum_{\alpha \in Q} \left( \sum_{\substack{\phi:Q'\rightarrow \mathbb{Z}_{\geq 0} \\ \sum_{\beta \in Q\setminus \{0\}} \phi(\beta)\beta = \alpha } } z^{\sum_{\beta \in Q\setminus \{0\}}\phi(\beta)}  \right) x^{\alpha }\\
    &= \frac{1}{1-z} \sum_{\mu} \left( \sum_{\substack{\phi:Q'\rightarrow \mathbb{Z}_{\geq 0} \\ \sum_{\beta \in Q\setminus \{0\}} \phi(\beta)\beta = \mu } } z^{\sum_{\beta \in Q\setminus \{0\}}\phi(\beta)}  \right) m_{\mu}(X).\\
\end{align*}

Define the series
$$f_{\mu}(z):= \sum_{\substack{\phi:Q'\rightarrow \mathbb{Z}_{\geq 0} \\ \sum_{\beta \in Q\setminus \{0\}} \phi(\beta)\beta = \mu } } z^{\sum_{\beta \in Q\setminus \{0\}}\phi(\beta)}.$$ Note that $f_{\mu}(z)$ is a polynomial with degree $|\mu|.$ The leading coefficient of $z^{|\mu|}$ is always $1$ corresponding to the vector partition $\mu= \mu_1(1,0\ldots) + \mu_2 (0,1,0,\ldots) + \ldots + \mu_{\ell(\mu)}(0,\ldots,0,1,0,\ldots)$. Now we see 
$$\Omega\left(z\Omega(X) \right) = \sum_{\mu} \frac{f_{\mu}(z)}{1-z} m_{\mu}(X).$$ This means that 
$$\langle h_{\mu}(X), \Omega\left(z\Omega(X) \right) \rangle = \frac{f_{\mu}(z)}{1-z}.$$ Putting this all together, 
\begin{align*}
    &R_{\lambda}(z) = \langle s_{\lambda}(X), \Omega\left(z \Omega(X) \right) \rangle\\
    &= \sum_{\mu\vdash |\lambda|} K^{(-1)}_{\lambda,\mu} \langle h_{\mu}(X), \Omega\left(z \Omega(X) \right) \rangle\\
    &= \sum_{\mu\vdash |\lambda|} K^{(-1)}_{\lambda,\mu} \frac{f_{\mu}(z)}{1-z}\\
    &= \frac{\sum_{\mu\vdash |\lambda|} K^{(-1)}_{\lambda,\mu} f_{\mu}(z)}{1-z}.\\
\end{align*}

It is clear that $\sum_{\mu\vdash n} K^{(-1)}_{\lambda,\mu} f_{\mu}(z)$ is a polynomial with degree at most $|\lambda|.$

\end{proof}

\end{document}
